Estimation of $(\mathbf{T}, \{\boldsymbol\theta_k\}, \boldsymbol\pi)$ proceeds by EM in log-space.
This section gives the shared scaffold (Section~\ref{sec:em_scaffold}), the family-specific M-step updates (Section~\ref{sec:mstep}), and the state-count selection protocol (Section~\ref{sec:k_selection}).

\subsection{Shared EM Scaffold}
\label{sec:em_scaffold}

All three emission families share the same log-space forward-backward
recursions and the same quantile-based initialization \cite{bilmes1998gentle}:
observations are sorted, partitioned into $K$ equal-sized chunks, and each
state's initial location and scale are seeded from the corresponding chunk;
$T_{ij}^{(0)} = \pi_k^{(0)} = 1/K$.
EM iterates to tolerance $|\mathcal{L}^{(n)} - \mathcal{L}^{(n-1)}| < 10^{-4}$
with $\texttt{max\_iter} = 60$; convergence is typically reached in
20--40 iterations.
Appendix~\ref{sec:baum_welch_supp} gives the forward-backward recursion,
the posterior quantities $\gamma_t(k)$ and $\xi_t(i, j)$, and the
family-specific M-step updates.

\subsection{Family-Specific M-Step Updates}
\label{sec:mstep}

\paragraph{CHMM-N (Gaussian).}
The classical Baum-Welch M-step \cite{baum1970maximization} is closed form:
$\mu_k \leftarrow \sum_t \gamma_t(k) O_t / \sum_t \gamma_t(k)$ and
$\sigma_k^2 \leftarrow \sum_t \gamma_t(k) (O_t - \mu_k)^2 / \sum_t \gamma_t(k)$.

\paragraph{CHMM-t (Student-t via ECM).}
We adopt the ECM formulation of \citet{peel2000robust} and \citet{liu1995ml}, the same family of heavy-tailed innovation choices studied outside the HMM framework by \citet{abantovalle2017svm} for stochastic-volatility-in-mean models:
the latent precision
\begin{equation}
    u_{t,k} = \frac{\nu_k + 1}{\nu_k + ((O_t - \mu_k)/\sigma_k)^2}
    \label{eq:tprecision}
\end{equation}
is treated as an augmented E-step quantity, and the M-step updates
\begin{equation}
    \mu_k \leftarrow \frac{\sum_t \gamma_t(k) u_{t,k} O_t}{\sum_t \gamma_t(k) u_{t,k}},
    \qquad
    \sigma_k^2 \leftarrow \frac{\sum_t \gamma_t(k) u_{t,k} (O_t - \mu_k)^2}{\sum_t \gamma_t(k)}
    \label{eq:tupdate}
\end{equation}
have closed form conditional on $\nu_k$.
The degrees of freedom $\nu_k$ are updated by a golden-section search
maximizing the per-state $Q$-function
$Q_k(\nu) = \sum_t \gamma_t(k) \log t_\nu(O_t; \mu_k, \sigma_k)$
over the bracket $(\nu_{\min}, \nu_{\max}) = (2.1, 50)$.
This M-step preserves the monotonicity of EM while avoiding the joint
saddle-point issues of simultaneous $(\mu_k, \sigma_k, \nu_k)$ maximization.

\paragraph{CHMM-L (Laplace).}
The weighted maximum-likelihood estimators are in closed form:
$\mu_k$ is the posterior-weighted median of the observations
(weights $\gamma_t(k)$, computed by cumulative-weight bracketing with
linear interpolation between neighboring order statistics), and
$b_k \leftarrow \sum_t \gamma_t(k) |O_t - \mu_k| / \sum_t \gamma_t(k)$.

\subsection{State-Count Selection}
\label{sec:k_selection}

We sweep $K \in \{3, 6, 9, 12, 15, 18, 21\}$ for CHMM-N and confirm on
CHMM-t / CHMM-L that the selected $K$ sits inside the IC-minimizing band.
We compute the four information criteria used for HMM state selection
in \citet{nguyen2018hidden},
\begin{equation}
    \text{AIC} = -2\mathcal{L} + 2p,\quad
    \text{BIC} = -2\mathcal{L} + p \ln T,\quad
    \text{HQC} = -2\mathcal{L} + 2p \ln(\ln T),\quad
    \text{CAIC} = -2\mathcal{L} + p(\ln T + 1),
    \label{eq:information_criteria}
\end{equation}
with $p = K^2 + 2K - 1$ for CHMM-N and CHMM-L, and $p = K^2 + 3K - 1$
for CHMM-t (one additional per-state degree-of-freedom parameter);
$\mathcal{L}$ is the converged log-likelihood.
The operating point is the $K$ that minimizes the average information-criterion
rank and simultaneously wins or ties every IS and OoS
distributional pass-rate metric.
